\documentclass[10pt,a4paper]{article}
\usepackage[T1]{fontenc}
\usepackage[utf8]{inputenc}
\usepackage{lmodern}
\usepackage[a4paper,left=0.55in,right=0.55in,top=0.42in,bottom=0.42in]{geometry}
\usepackage{enumitem}
\usepackage{xcolor}
\definecolor{linkblue}{HTML}{1F5F8B}
\usepackage[colorlinks=true,linkcolor=black,urlcolor=linkblue,citecolor=black]{hyperref}
\usepackage{microtype}
\usepackage{tabularx}

\pagestyle{empty}
\setlength{\parindent}{0pt}
\setlength{\parskip}{0pt}
\setlength{\tabcolsep}{0pt}
\emergencystretch=1em

\newcommand{\cvheader}[2]{%
  \begin{center}
    {\fontsize{17}{18}\selectfont\bfseries Rody Vilchez}\par
    \vspace{1pt}
    {\normalsize Applied ML Engineer --- Retrieval $\cdot$ Agent Systems $\cdot$ Evaluation $\cdot$ Urban ML}\par
    \vspace{1pt}
    {\small
      \href{mailto:rody@rosewt.dev}{rody@rosewt.dev} \;|\;
      +51 987 082 126 \;|\;
      \href{https://rosewt.dev}{rosewt.dev} \;|\;
      \href{https://github.com/R0SEWT}{GitHub} \;|\;
      \href{https://www.linkedin.com/in/r0sewt/}{LinkedIn}}
  \end{center}
  \vspace{3pt}
  #1\par
  \vspace{5pt}
}

\newcommand{\cvsection}[1]{%
  \vspace{6pt}%
  {\fontsize{12}{13}\selectfont\bfseries #1}\par
  \vspace{1pt}\hrule height 0.35pt\vspace{4pt}%
}

\newcommand{\role}[4]{%
  \begin{tabular*}{\textwidth}{@{\extracolsep{\fill}}lr}
    \textbf{#1} & \textbf{#2} \\
    \textit{#3} & \textit{#4}
  \end{tabular*}\par
  \vspace{-2pt}%
}

\newcommand{\project}[2]{%
  \begin{tabular*}{\textwidth}{@{\extracolsep{\fill}}lr}
    \textbf{#1} & \textit{#2}
  \end{tabular*}\par
  \vspace{-2pt}%
}

\newenvironment{cvitems}{%
  \begin{itemize}[leftmargin=1.25em,labelsep=0.45em,itemsep=2pt,topsep=2.5pt,parsep=0pt,partopsep=0pt]
}{%
  \end{itemize}\vspace{-1pt}
}

\newcommand{\lead}[1]{\textbf{#1} ---}


\hypersetup{
  pdftitle={Rody Vilchez - CV},
  pdfauthor={Rody Vilchez}
}

\begin{document}

\cvheader{%
Ingeniero de ML aplicado en la intersección de sistemas agénticos, retrieval y evaluación bajo datos imperfectos. En
CIP--CGIAR construyo plataformas multiagente en producción, GraphRAG multilingüe e infraestructura de evaluación;
en investigación, estudio fallos de medición en urban ML y soy coautor de una arquitectura multimodal publicada en Springer CCIS.
}{}

\cvsection{Experiencia}

\role{AI Intern}{International Potato Center (CIP--CGIAR)}{Lima, Perú}{Oct 2025 -- Presente}
\begin{cvitems}
  \item \lead{Plataforma multiagente en Microsoft Teams} Arquitecté y llevé a producción sobre Copilot Studio un sistema
  institucional para soporte TI y servicios corporativos, con delegación entre dominios y escalamiento a tickets
  prellenados desde el contexto conversacional, sujetos a revisión humana.
  \item \lead{Suite de evaluación agent-based} Construí una suite donde agentes evaluadores ejecutan escenarios derivados
  de registros operativos; combina regresión determinista, LLM-as-judge agnóstico de proveedor y simulación por perfiles
  para atribuir fallos de ruteo, delegación y calidad de respuesta.
  \item \lead{GraphRAG multilingüe} Construí end-to-end la capa de document intelligence que alimenta un GraphRAG
  agrícola en cinco idiomas sobre corpus con OCR ruidoso y layouts irregulares: parsing, chunking, extracción estructurada
  de metadatos, embeddings y almacenamiento vectorial.
\end{cvitems}

\role{QA Trainee}{Visma LATAM}{Lima, Perú}{Dic 2024 -- Oct 2025}
\begin{cvitems}
  \item \lead{Síntesis de pruebas con LLM} Construí un agente que transforma especificaciones funcionales en escenarios
  E2E ejecutables.
  \item \lead{Automatización DOM-aware} Desarrollé extracción de selectores y estado del DOM, integrando suites Cypress
  generadas en Jenkins como cobertura de regresión para flujos críticos.
\end{cvitems}

\cvsection{Investigación y sistemas seleccionados}

\project{Infelix --- Riesgo delictivo latente bajo sesgo de medición}{Investigación de tesis $\cdot$ UPC $\cdot$ 2025--2026}
\begin{cvitems}
  \item \lead{Modelado del delito no denunciado} Construí un modelo bayesiano sobre encuestas de victimización y un panel
  H3 de 11 fuentes geoespaciales (Sentinel-2, VIIRS, OSM, WorldPop y Street View) para estimar riesgo delictivo en 43
  distritos de Lima--Callao; los datos censales fueron la única fuente que mejoró consistentemente las predicciones
  dentro de cada distrito.
  \item \lead{Transferencia entre ciudades} Las features aprendidas en Lima igualaron modelos con un año de datos locales
  en Arequipa y se mantuvieron competitivas en Cusco y Piura; las pruebas mostraron que la calidad del registro, no solo
  la arquitectura, determina cuándo un sistema puede evaluarse y transferirse.
\end{cvitems}

\project{Multimodal Sign Language Model}{Springer CCIS 2895, 2026}
\begin{cvitems}
  \item Coautor de una arquitectura multimodal publicada, sin glosas, que usa queries latentes y cross-attention para proyectar
  keypoints 2D al espacio de embeddings de un modelo de lenguaje preentrenado, evaluada en corpus peruanos y
  argentinos. DOI: \href{https://doi.org/10.1007/978-3-032-20322-9_23}{10.1007/978-3-032-20322-9\_23}
\end{cvitems}

\project{Wachi --- Ruteo peatonal sensible al riesgo}{El Comercio Lab}
\begin{cvitems}
  \item Desarrollé un sistema geoespacial que reordena rutas peatonales por riesgo espaciotemporal en Lima mediante
  superficies H3, decaimiento temporal y seis franjas horarias para una experiencia interactiva de periodismo de datos.
\end{cvitems}

\cvsection{Contribuciones open source}
\begin{cvitems}
  \item \textbf{Microsoft Copilot Studio} --- Diagnostiqué fallos de sincronización de knowledge sources en subagentes y
  validé correcciones de prerelease en la extensión oficial de VS Code.
  \item \textbf{Gemini CLI} --- Corregí un esquema de UI y añadí cobertura de regresión al CLI agéntico open source de Google.
  \item \textbf{scikit-learn} --- Mejoré la documentación técnica de HDBSCAN.
  \item \textbf{beads (bd)} --- Implementé prefix-routed write-through y el test de regresión que desbloqueó la consolidación
  de la implementación upstream.
\end{cvitems}

\cvsection{Educación y reconocimientos}
\begin{tabular*}{\textwidth}{@{\extracolsep{\fill}}lr}
  \textbf{Ciencias de la Computación} & \textbf{Universidad Peruana de Ciencias Aplicadas (UPC)}
\end{tabular*}\par
\textit{Graduación prevista: fines de 2026}\par
Segundo puesto, DataFest BCP $\times$ ESAN 2025 $\cdot$ Beca completa SALA 2026

\cvsection{Foco técnico}
\textbf{Retrieval y agentes:} LlamaIndex, Qdrant, Copilot Studio, Power Automate, orquestación multiagente\par
\textbf{Evaluación e investigación:} Python, PyTorch, scikit-learn, LLM-as-judge, diseño experimental, robustez e incertidumbre\par
\textbf{Datos y backend:} FastAPI, H3, DuckDB, PostgreSQL, ETL, APIs REST, Docker, GitHub Actions, Linux\par
\textbf{Idiomas:} Español (nativo) $\cdot$ Inglés (intermedio)

\end{document}
